% ##################################################################
% CAPÍTULO I: INTRODUCCIÓN
% ##################################################################
\section{Introducci{\'o}n}
En la regi{\'o}n Cusco, donde m{\'a}s de la mitad de la poblaci{\'o}n tiene el quechua como lengua
materna, la atenci{\'o}n en salud opera de manera predominante en espa{\~n}ol. Esta asimetr{\'i}a
levanta una barrera comunicativa entre un personal de salud mayoritariamente castellanohablante y una
poblaci{\'o}n usuaria quechuahablante, barrera que ---agravada por la baja alfabetizaci{\'o}n en las
zonas rurales andinas--- compromete aspectos sensibles del acto m{\'e}dico, desde la recogida de los
s{\'i}ntomas hasta la correcta comprensi{\'o}n de las indicaciones.

En paralelo, el procesamiento del lenguaje natural ha alcanzado un punto en el que es posible
\emph{adaptar} grandes modelos multiling{\"u}es preentrenados a lenguas de bajo recurso como el
quechua, sin necesidad de entrenarlos desde cero. Sobre esa base, el presente plan propone desarrollar
y evaluar un sistema \textbf{bidireccional de traducci{\'o}n de voz Quechua--Espa{\~n}ol} para
consultas m{\'e}dicas, que encadena reconocimiento de voz, traducci{\'o}n autom{\'a}tica y s{\'i}ntesis
de voz, con salida hablada en ambos sentidos del di{\'a}logo para incluir tambi{\'e}n a los usuarios
con baja o nula alfabetizaci{\'o}n.

Este primer cap{\'i}tulo sit{\'u}a el problema en sus niveles internacional, nacional y regional,
delimita el vac{\'i}o de investigaci{\'o}n y el problema espec{\'i}fico, y formula la propuesta con sus
objetivos, su justificaci{\'o}n y sus alcances. Las hip{\'o}tesis, por su car{\'a}cter contrastable, se
enuncian en el Cap{\'i}tulo~III junto con el dise{\~n}o metodol{\'o}gico que permite someterlas a
prueba.

\subsection{Contexto y problema general}
\obsdos{El contexto de esta investigaci{\'o}n se aborda en tres niveles ---internacional, nacional y
regional--- que, en conjunto, muestran c{\'o}mo una barrera ling{\"u}{\'i}stica hist{\'o}rica se
combina hoy con una oportunidad tecnol{\'o}gica inexistente hace pocos a{\~n}os.}

\subsubsection{\obsdos{Contexto internacional}}
\begin{bloqueObsDos}
\noindent De las miles de lenguas que se hablan en el mundo, solo una peque{\~n}a fracci{\'o}n cuenta
con recursos digitales suficientes para sostener tecnolog{\'i}as del lenguaje; la gran mayor{\'i}a
son \emph{lenguas de bajo recurso}, con escasa o nula presencia en los sistemas de voz y
traducci{\'o}n autom{\'a}tica \parencite{joshi2020diversity}. Esta asimetr{\'i}a digital reproduce, en
el plano tecnol{\'o}gico, las desigualdades que ya enfrentan sus hablantes en el acceso a servicios
b{\'a}sicos.

\noindent En el {\'a}mbito
sanitario, la comunicaci{\'o}n entre un profesional de salud y un paciente que no comparten la misma
lengua es un problema reconocido a escala global: la barrera idiom{\'a}tica compromete la
anamnesis, el diagn{\'o}stico, el consentimiento informado y la adherencia al tratamiento
\parencite{alshamsi2020barreras}, y por ello el derecho de la persona a ser atendida en su lengua
materna se considera un componente de la calidad y la seguridad del acto m{\'e}dico.

\noindent En paralelo, el procesamiento del lenguaje natural ha experimentado un cambio de
paradigma que abre una ventana de oportunidad para estas lenguas. La estrategia dominante ya no es
entrenar modelos desde cero ---inviable con pocos datos---, sino adaptar (\emph{fine-tuning})
grandes modelos multiling{\"u}es preentrenados que incorporan cientos de idiomas: en la voz, MMS
\parencite{pratap2024mms} y Whisper \parencite{radford2023whisper}; en la traducci{\'o}n, NLLB
\parencite{costajussa2022nllb}.

\noindent Competencias internacionales como AmericasNLP
\parencite{mager2021americasnlp} y las tareas de traducci{\'o}n de voz de IWSLT
\parencite{Ortega2025260} han impulsado de forma sostenida el trabajo sobre lenguas ind{\'i}genas
americanas. En este panorama, el quechua ---la lengua originaria con mayor n{\'u}mero de hablantes
de Am{\'e}rica, extendida por varios pa{\'i}ses andinos--- concentra buena parte de ese esfuerzo,
aunque todav{\'i}a de manera fragmentaria y por componentes aislados.

\noindent En el terreno sanitario, la pr{\'a}ctica internacional recomendada frente a la barrera
idiom{\'a}tica es la mediaci{\'o}n ling{\"u}{\'i}stica ---int{\'e}rpretes presenciales o
telef{\'o}nicos--- \parencite{alshamsi2020barreras}; sin embargo, su disponibilidad es limitada y
pr{\'a}cticamente nula en entornos rurales y de atenci{\'o}n primaria, donde ocurre el primer
contacto con el sistema de salud.

\noindent Por eso
las tecnolog{\'i}as del habla ---reconocimiento, traducci{\'o}n y s{\'i}ntesis de voz--- se perfilan
a escala global como un apoyo emergente a la comunicaci{\'o}n cl{\'i}nica, con la ventaja decisiva de
operar por \emph{voz}: a diferencia de las soluciones basadas en texto, no exigen alfabetizaci{\'o}n
y resultan utilizables por poblaciones con baja escolaridad.

\noindent Esta orientaci{\'o}n se inscribe en un
esfuerzo m{\'a}s amplio de \emph{inclusi{\'o}n digital} y de preservaci{\'o}n de las lenguas
originarias, hoy reconocido como un objetivo de equidad tecnol{\'o}gica: dotar de presencia digital a
una lengua es tambi{\'e}n una forma de sostener los derechos de sus hablantes.
\end{bloqueObsDos}

\subsubsection{\obsdos{Contexto nacional}}
El Per{\'u} es un pa{\'i}s plurilingüe en el que el quechua es la lengua originaria con mayor
n{\'u}mero de hablantes. Los Censos Nacionales 2017 registran a millones de peruanos que declaran
el quechua como lengua materna, con una marcada concentraci{\'o}n en la sierra sur
\parencite{inei2017}.
\begin{bloqueObsDos}
\noindent El Estado reconoce el car{\'a}cter multiling{\"u}e del pa{\'i}s y el derecho de los
ciudadanos a relacionarse con las instituciones p{\'u}blicas en su lengua materna; sin embargo, en
la pr{\'a}ctica los servicios esenciales ---la salud entre ellos--- operan de manera predominante en
espa{\~n}ol, lo que condiciona su acceso efectivo para la poblaci{\'o}n quechuahablante.

\noindent A esta
brecha ling{\"u}{\'i}stica se suma una brecha educativa: cerca del 22.8\,\% de las mujeres rurales no
sabe leer ni escribir \parencite{inei2023}, con programas de alfabetizaci{\'o}n a{\'u}n limitados en
las zonas andinas \parencite{dispurse2023}.

\noindent Esta doble barrera ---de idioma y de
alfabetizaci{\'o}n--- implica que una soluci{\'o}n tecnol{\'o}gica pertinente para el pa{\'i}s no
puede basarse solo en texto, sino que debe operar principalmente por \emph{voz}. En el plano
institucional ya existen se{\~n}ales de respuesta: EsSalud reactiv{\'o} en 2023 una central
telef{\'o}nica de atenci{\'o}n en quechua \parencite{essalud2023}, y la comunidad acad{\'e}mica
peruana ha producido recursos y modelos iniciales para el quechua en reconocimiento, traducci{\'o}n
y s{\'i}ntesis de voz, aunque ninguno orientado a la consulta m{\'e}dica.

\noindent El sistema de salud peruano incorpora, al menos en el plano normativo, un enfoque
intercultural que reconoce la pertinencia ling{\"u}{\'i}stica de la atenci{\'o}n
\parencite{minsa2016intercultural}; no obstante, su
traducci{\'o}n en herramientas operativas dentro del consultorio sigue siendo escasa. La atenci{\'o}n
primaria se presta en una red amplia y dispersa de establecimientos ---postas y centros de salud---
situados en zonas andinas donde la poblaci{\'o}n usuaria es mayoritariamente quechuahablante, lo que
convierte la distancia idiom{\'a}tica en un problema cotidiano y estructural, no excepcional.

\noindent En ese
contexto, la investigaci{\'o}n nacional en procesamiento del lenguaje para el quechua ha crecido de
forma sostenida en la {\'u}ltima d{\'e}cada, pero se ha concentrado en resolver cada componente por
separado ---reconocimiento, traducci{\'o}n o s{\'i}ntesis--- y sobre dominios generales, sin integrar
un sistema completo ni abordar el vocabulario propio de la atenci{\'o}n en salud.
\end{bloqueObsDos}

\subsubsection{\obsdos{Contexto regional (Cusco)}}
En la regi{\'o}n Cusco la presencia del quechua es particularmente intensa: m{\'a}s del
55\,\% de la poblaci{\'o}n tiene el quechua como lengua materna \parencite{inei2017}. Esta
realidad ling{\"u}{\'i}stica condiciona el acceso efectivo a los servicios p{\'u}blicos
esenciales ---la atenci{\'o}n en salud entre ellos--- cuando operan de manera predominante en
espa{\~n}ol.

\begin{bloqueObsDos}
\noindent La red de salud de Cusco atiende a una poblaci{\'o}n extensa y geogr{\'a}ficamente
dispersa, con numerosas comunidades rurales y altoandinas donde el quechua es la lengua de uso
cotidiano y el espa{\~n}ol tiene presencia limitada.

\noindent A esta realidad se suma una particularidad
ling{\"u}{\'i}stica decisiva para la propuesta: el quechua de Cusco corresponde a la variante Collao
(\sig{quz}), distinta de la de Ayacucho (\sig{quy}) tanto en l{\'e}xico como en rasgos
fonol{\'o}gicos; esta diferencia dialectal ---no menor para un sistema autom{\'a}tico--- es
justamente la que los modelos de traducci{\'o}n disponibles a{\'u}n no cubren para Cusco.

\noindent Por
a{\~n}adidura, la consulta m{\'e}dica no es solo un intercambio de informaci{\'o}n, sino un encuentro
cultural en el que la confianza y la comprensi{\'o}n mutua condicionan la adherencia al tratamiento;
cuando esa comunicaci{\'o}n falla, el riesgo cl{\'i}nico recae sobre la poblaci{\'o}n m{\'a}s
vulnerable.
\end{bloqueObsDos}

\subsubsection{\obsdos{Planteamiento del problema}}
\obsdos{Los tres niveles de contexto ---internacional, nacional y regional--- convergen en un
problema concreto y localizado: la comunicaci{\'o}n cl{\'i}nica en la consulta cara a cara de los
establecimientos de salud de Cusco, donde el idioma se vuelve un determinante directo de la
calidad y la seguridad de la atenci{\'o}n.}

En el {\'a}mbito sanitario, la distancia idiom{\'a}tica entre un personal de salud mayoritariamente
castellanohablante y una poblaci{\'o}n usuaria quechuahablante origina una barrera comunicativa con
consecuencias cl{\'i}nicas directas.

EsSalud reconoce que cerca de 70\,000 de sus asegurados en
Cusco son quechuahablantes y reactiv{\'o} en 2023 la central telef{\'o}nica en quechua
\parencite{essalud2023}; sin embargo, \obsdos{esa medida ---de alcance telef{\'o}nico y
unidireccional--- no alcanza} la interacci{\'o}n cara a cara del
consultorio, donde se recogen s{\'i}ntomas, se formula el diagn{\'o}stico y se entregan las
indicaciones.

La evidencia confirma la gravedad del problema: \textcite{begazo2026} documenta
c{\'o}mo la desconexi{\'o}n idiom{\'a}tica vulnera la relaci{\'o}n terap{\'e}utica en hospitales de
Cusco, y \textcite{cuaila2022} halla una relaci{\'o}n entre el conocimiento del quechua del personal
y la calidad de atenci{\'o}n percibida. A ello se suma una barrera de alfabetizaci{\'o}n: cerca del
22.8\,\% de las mujeres rurales no sabe leer ni escribir \parencite{inei2023}, lo que obliga a una
soluci{\'o}n basada en \emph{voz} y no solo en texto (Figura~\ref{fig:barreras}).

\obsdos{\noindent Las consecuencias de esta barrera no son menores: una anamnesis incompleta, un
diagn{\'o}stico tard{\'i}o o una indicaci{\'o}n mal comprendida comprometen la seguridad del acto
m{\'e}dico y la adherencia al tratamiento, y recaen sobre la poblaci{\'o}n m{\'a}s vulnerable. Pese a
ello, las respuestas institucionales vigentes son parciales ---telef{\'o}nicas, unidireccionales o
dependientes de la disponibilidad de personal biling{\"u}e--- y no existe una herramienta que asista
la traducci{\'o}n por voz, en ambos sentidos, durante la propia consulta.}

\begin{bloqueObsDos}
\noindent En s{\'i}ntesis, el problema general que motiva este plan es la \textbf{barrera
comunicativa entre el personal de salud castellanohablante y los pacientes quechuahablantes de
Cusco} durante la consulta cara a cara, agravada por la baja alfabetizaci{\'o}n y no resuelta por
las medidas vigentes. Existiendo hoy modelos multiling{\"u}es preentrenados que cubren parcialmente
el quechua, la pregunta es si pueden adaptarse e integrarse en un sistema de traducci{\'o}n por voz,
bidireccional, para asistir esa interacci{\'o}n cl{\'i}nica.
\end{bloqueObsDos}

\begin{figure}[H]
\centering
\caption{Magnitud de la problem{\'a}tica en Cusco.}
\label{fig:barreras}
\begin{tikzpicture}
\begin{axis}[
    ybar, width=12cm, height=6cm, bar width=28pt,
    ymin=0, ymax=100, ylabel={Porcentaje (\%)},
    symbolic x coords={Lengua materna,Asegurados EsSalud,Analfab. mujer rural},
    xtick={Lengua materna,Asegurados EsSalud,Analfab. mujer rural},
    nodes near coords, nodes near coords style={font=\small},
    x tick label style={align=center,font=\small,text width=2.5cm},
    enlarge x limits=0.25, ymajorgrids=true, grid style=dashed,
]
\addplot[fill=g4, draw=g5, bar shift=0pt] coordinates {(Lengua materna,55)};
\addplot[fill=g8, draw=g7, bar shift=0pt] coordinates {(Asegurados EsSalud,15)};
\addplot[fill=g12, draw=g11, bar shift=0pt] coordinates {(Analfab. mujer rural,22.8)};
\end{axis}
\end{tikzpicture}
\notaAPA{Elaboraci{\'o}n propia con datos de \textcite{inei2017}, \textcite{essalud2023} e \textcite{inei2023}.}
\end{figure}

\subsection{Tecnolog{\'i}a o m{\'e}todo}
El estado del arte en procesamiento del lenguaje natural ofrece hoy una oportunidad que no
exist{\'i}a hace pocos a{\~n}os. La estrategia dominante para lenguas de bajo recurso ya no consiste
en entrenar modelos desde cero, sino en adaptar (\emph{fine-tuning}) grandes modelos multiling{\"u}es
preentrenados.

En la voz, el proyecto MMS \parencite{pratap2024mms} cubre el quechua de Cusco
(\sig{quz}) en reconocimiento y s{\'i}ntesis; Whisper \parencite{radford2023whisper} declara una
cobertura marginal del quechua y se reserva como l{\'i}nea base. En traducci{\'o}n, NLLB
\parencite{costajussa2022nllb} incluye el quechua de Ayacucho (\sig{quy}) y no el de Cusco, por lo
que el par debe adaptarse a \sig{quz} ---esa brecha dialectal es el n{\'u}cleo del aporte
t{\'e}cnico---.

Existen adem{\'a}s antecedentes de traducci{\'o}n neural para esta lengua
polisint{\'e}tica \parencite{ortega2020quechua,mager2021americasnlp} y recursos como el corpus
\emph{Huqariq} \parencite{zevallos2022huqariq}.

\subsection{Vac{\'i}o de investigaci{\'o}n}
Persiste una desconexi{\'o}n entre ambos mundos: el problema social est{\'a} bien documentado y la
tecnolog{\'i}a base existe, pero nadie ha unido los dos extremos. No se dispone de un corpus
cl{\'i}nico Quechua--Espa{\~n}ol ni de una evaluaci{\'o}n emp{\'i}rica del desempe{\~n}o de estos
modelos sobre el vocabulario propio de la consulta m{\'e}dica (triaje, s{\'i}ntomas e indicaciones),
y NLLB a{\'u}n no cubre la variante de Cusco (\sig{quz}). Cerrar esa intersecci{\'o}n es el aporte
de este proyecto.

\subsection{Formulaci{\'o}n del problema}
La formulaci{\'o}n del problema de esta investigaci{\'o}n se concentra en la comunicaci{\'o}n
m{\'e}dico-paciente mediada por un sistema bidireccional de traducci{\'o}n de voz Quechua--Espa{\~n}ol
basado en modelos preentrenados, y en la medida en que su ajuste al dominio cl{\'i}nico mejora la
precisi{\'o}n y la comprensi{\'o}n frente a la ausencia de apoyo tecnol{\'o}gico.

\subsubsection{Problema general}
{\textquestiondown}\obsdos{C{\'o}mo \emph{desarrollar un sistema bidireccional de traducci{\'o}n de voz
Quechua--Espa{\~n}ol} \emph{mediante arquitecturas de transformadores multimodales} (ajuste fino de los
modelos preentrenados MMS y NLLB) \emph{para la comunicaci{\'o}n m{\'e}dico-paciente en las consultas m{\'e}dicas en Cusco}?}

\subsubsection{Problemas espec{\'i}ficos}
\begin{enumerate}
  \item \obsdos{\textbf{(PE1 --- Diagn{\'o}stico).} {\textquestiondown}Cu{\'a}les son las barreras de comunicaci{\'o}n entre el personal de salud y los pacientes quechuahablantes, y qu{\'e} debe contener el corpus cl{\'i}nico Quechua--Espa{\~n}ol (triaje, s{\'i}ntomas e indicaciones) para poder ajustar los modelos? \textbf{Aporte:} Permite construir el corpus cl{\'i}nico etiquetado; sin esos datos no es posible ajustar los modelos.}
  \item \obsdos{\textbf{(PE2 --- Construcci{\'o}n).} {\textquestiondown}C{\'o}mo construir el sistema de traducci{\'o}n por voz, ajustando (\emph{fine-tuning}) los modelos preentrenados MMS y NLLB ---del quechua de Ayacucho (\sig{quy}) al de Cusco (\sig{quz})--- e integr{\'a}ndolos en un \emph{pipeline} ASR+MT+TTS? \textbf{Aporte:} Es el n{\'u}cleo de la soluci{\'o}n: el sistema con el modelo de traducci{\'o}n adaptado al dominio cl{\'i}nico y a la variante de Cusco.}
  \item \obsdos{\textbf{(PE3 --- Validaci{\'o}n).} {\textquestiondown}En qu{\'e} medida el sistema mejora la precisi{\'o}n de la traducci{\'o}n (WER, BLEU, chrF), la calidad de la voz (MOS) y la comprensi{\'o}n entre el m{\'e}dico y el paciente, frente a no contar con apoyo tecnol{\'o}gico? \textbf{Aporte:} Demuestra si el sistema resuelve el problema general.}
\end{enumerate}

\subsection{Propuesta de soluci{\'o}n}
Se propone desarrollar y evaluar un sistema bidireccional de traducci{\'o}n de voz
Quechua--Espa{\~n}ol para consultas m{\'e}dicas, que encadena reconocimiento de voz (MMS),
traducci{\'o}n autom{\'a}tica (NLLB adaptado de \sig{quy} a \sig{quz}) y s{\'i}ntesis de voz
(MMS-TTS) mediante el ajuste de modelos preentrenados de c{\'o}digo abierto, con salida por voz en
ambos sentidos del di{\'a}logo.

\subsection{Objetivos}

\subsubsection{Objetivo general}
\nuevo{Desarrollar un sistema bidireccional de traducci{\'o}n de voz Quechua--Espa{\~n}ol para la
comunicaci{\'o}n m{\'e}dico-paciente, mediante arquitecturas de transformadores multimodales (ajuste fino
de los modelos preentrenados MMS y NLLB), en las consultas m{\'e}dicas en Cusco.}

\subsubsection{Objetivos espec{\'i}ficos}
\begin{enumerate}
  \item \obsdos{\textbf{(OE1 --- Diagn{\'o}stico).} Caracterizar las barreras de comunicaci{\'o}n entre el personal de salud y los pacientes quechuahablantes y determinar qu{\'e} debe contener el corpus cl{\'i}nico Quechua--Espa{\~n}ol (triaje, s{\'i}ntomas e indicaciones) para poder ajustar los modelos.}
  \item \obsdos{\textbf{(OE2 --- Construcci{\'o}n).} Construir el sistema de traducci{\'o}n por voz, ajustando (\emph{fine-tuning}) los modelos preentrenados MMS y NLLB ---del quechua de Ayacucho (\sig{quy}) al de Cusco (\sig{quz})--- e integr{\'a}ndolos en un \emph{pipeline} ASR+MT+TTS con salida por voz y texto.}
  \item \obsdos{\textbf{(OE3 --- Validaci{\'o}n).} Evaluar en qu{\'e} medida el sistema mejora la precisi{\'o}n de la traducci{\'o}n (WER, BLEU, chrF), la calidad de la voz (MOS) y la comprensi{\'o}n entre el m{\'e}dico y el paciente, frente a no contar con apoyo tecnol{\'o}gico.}
\end{enumerate}

\subsection{Justificaci{\'o}n}
\noindent La investigaci{\'o}n se justifica en cuatro dimensiones:
\begin{itemize}\setlength{\itemsep}{1pt}
  \item \textbf{Te{\'o}rica:} Aporta evidencia emp{\'i}rica del desempe{\~n}o de modelos
  preentrenados en una lengua de bajo recurso y morfolog{\'i}a compleja.
  \item \textbf{Pr{\'a}ctica:} Entrega un prototipo piloteable en un centro de salud.
  \item \textbf{Social:} Ataca una inequidad documentada en el acceso a la salud
  \parencite{begazo2026}, en l{\'i}nea con el derecho de los pacientes a ser atendidos en su
  lengua materna \parencite{essalud2023}.
  \item \textbf{Metodol{\'o}gica:} Establece un proceso replicable (corpus + ajuste + evaluaci{\'o}n)
  aplicable a otras variantes del quechua o a otros dominios.
\end{itemize}
\noindent Entre las \textbf{contribuciones originales esperadas} destacan:
\begin{itemize}\setlength{\itemsep}{1pt}
  \item El \textbf{primer corpus de dominio m{\'e}dico} Quechua--Espa{\~n}ol para tecnolog{\'i}as del habla.
  \item La \textbf{adaptaci{\'o}n de NLLB de \sig{quy} a \sig{quz}} en el dominio cl{\'i}nico
  (aporte t{\'e}cnico central).
  \item Un \textbf{sistema integrado ASR+MT+TTS} con salida por voz para usuarios con baja alfabetizaci{\'o}n.
\end{itemize}
\noindent En conjunto, el proyecto contribuye a la equidad en el acceso a la salud de la
poblaci{\'o}n quechuahablante de Cusco, a la seguridad del acto m{\'e}dico y a la revitalizaci{\'o}n
y el prestigio digital del quechua.

\subsubsection{Alcances y limitaciones}
\noindent\textbf{Alcances.}
\begin{itemize}
  \item Par ling{\"u}{\'i}stico Quechua de Cusco (\sig{quz}) -- Espa{\~n}ol, en ambos sentidos.
  \item Vocabulario acotado de consulta general (triaje, s{\'i}ntomas, indicaciones).
  \item Prototipo integrado ASR+MT+TTS con salida por voz y texto.
  \item Validaci{\'o}n con consultas simuladas en un centro de salud de referencia en Cusco.
\end{itemize}
\noindent\textbf{Limitaciones.}
\begin{itemize}
  \item \textbf{Temporal:} Ejecuci{\'o}n entre octubre 2026 y septiembre 2027 (12 meses).
  \item \textbf{{\'E}tica:} Se emplean di{\'a}logos simulados/actuados, sin datos reales de pacientes.
  \item Desempe{\~n}o condicionado por la escasez de datos de dominio cl{\'i}nico en quechua.
  \item Resultados circunscritos a la variante de Cusco; no generalizables sin ajuste a otras variantes.
\end{itemize}

\subsection{Organizaci{\'o}n de la tesis}
El documento se organiza en cinco cap{\'i}tulos:
\begin{itemize}\setlength{\itemsep}{1pt}
  \item \textbf{Cap{\'i}tulo~I --- Introducci{\'o}n:} contexto, tecnolog{\'i}a, vac{\'i}o de
  investigaci{\'o}n, problema, propuesta, objetivos y justificaci{\'o}n.
  \item \textbf{Cap{\'i}tulo~II --- Marco te{\'o}rico:} fundamentos te{\'o}ricos y estado del arte
  derivado de la revisi{\'o}n sistem{\'a}tica.
  \item \textbf{Cap{\'i}tulo~III --- Metodolog{\'i}a:} enfoque, dise{\~n}o, \obsdos{hip{\'o}tesis},
  variables, unidades de an{\'a}lisis, t{\'e}cnicas, procedimiento, protocolo de evaluaci{\'o}n y
  aspectos {\'e}ticos.
  \item \textbf{Cap{\'i}tulo~IV --- Factibilidad:} an{\'a}lisis de la factibilidad t{\'e}cnica,
  econ{\'o}mica y operativa.
  \item \textbf{Cap{\'i}tulo~V --- Aspectos administrativos:} cronograma, presupuesto y equipo.
\end{itemize}
\nuevo{\noindent La \textbf{matriz de consistencia}, incluida al cierre del Cap{\'i}tulo~III,
articula de forma integral el problema, los objetivos, las hip{\'o}tesis, las variables y la
metodolog{\'i}a.}
